\scanpage{89}
% loc 1–3
$A_1$ et $A_2$ \inserted{diagonales} engendrent $M_2$. Cela se réalise au moins loc. fpqc (en passant à la \deleted{solution} rev. des solutions de
% loc 4
\[ \lambda^2-\rho\lambda+1=0\quad\deleted{\unsure{\rho}} \]
% loc 5–8
d'autre part, à cause de $b^2-4c$ \unsure{inv.|\deleted{inv.}}, on peut trouver loc. fpqc (passant au rev. des solutions de $x^2-bx+c=0$), $u$ dans $A_1$ correspondant à
% loc 9
\[ \operatorname{Tr}u=b,\quad \det u=c \]
% loc 10
et item $v$ dans $A_2$
% loc 11
\[ Tv=b',\quad \det v=c'. \]
% loc 12
Soit \deleted{$t'$}
\[ \tau=\operatorname{Tr}uv \]
% loc 13–15
alors $\tau$ satisfait la relation analogue à $(11')$ avec $t$ remplacé par $\tau$. Or en tant qu'équation en $t$
% loc 16
\[ t^2-bb't+\left(b^2c'+b'^2c-4cc'-\frac{\delta\delta'}{\rho_0}\right) \]
% loc 17–19
(où $\delta=b^2-4c,\ \delta'=b'^2-4c'$), il y a \deleted{tout juste} (sur un corps) \deleted{au plus} 2 solutions, \deleted{soient} $t,t'$ (\unsure{évt.|év!} égales) satisfaisant
% loc 20
\[ t+t'=bb' \]
% loc 21–22, bottom
Si $t=\tau$, on a gagné. Sinon \inserted{$(\tau=t')$} on remplace $u$ par
% archive sheet number: 88
